
Language model alignment faces a fundamental tension between optimizing for aggregate human preferences---measured through pairwise comparisons and captured via Direct Preference Optimization \citep{Rafailov2023DPO}---and enabling user-specific customization through controllable attributes such as helpfulness, verbosity, and creativity \citep{Dong2023SteerLM}. Existing approaches treat these as separate objectives requiring sequential training: first optimize preferences via DPO, then fine-tune for attribute control. This sequential paradigm introduces catastrophic forgetting risks where the second objective degrades the first, forcing practitioners to choose between quality guarantees and user customization.

The conventional assumption is that DPO's implicit reward modeling and attribute conditioning's explicit user control are fundamentally incompatible for joint optimization. This assumption stems from observations that multi-task learning often exhibits negative transfer when task gradients conflict \citep{Yu2020PCGrad}. However, this incompatibility has never been empirically tested through gradient-level analysis of DPO and attribute objectives. Multi-task learning theory establishes that when task gradients maintain angles below 120°, joint optimization can succeed without catastrophic interference \citep{Navon2022NashMTL}. Whether DPO preference optimization and attribute conditioning satisfy this criterion remains unknown.

We hypothesize that these objectives are gradient-compatible, enabling joint training without destructive task conflict. To validate this, we implement a proof-of-concept joint training system with real-time gradient monitoring, measuring the angle between $\nabla L_DPO$ and $\nabla L_attr$ throughout training. Our experiments on GPT-2 XL (1.5B parameters) using the HH-RLHF preference dataset and OpenAssistant attribute annotations reveal a mean gradient angle of 78.5° (SD: 12.8°), well below the 120° interference threshold. Zero percent of gradient measurements across 100 training steps exceeded this threshold.

This gradient compatibility translates into practical capability. Both DPO loss and attribute loss decreased monotonically throughout training---DPO loss by 5.8\% and attribute loss by 21.3\%---confirming convergence without oscillation or objective degradation. The jointly trained model achieved 54.07\% preference win rate (approximately 94\% of standalone DPO baseline performance) while simultaneously achieving 65.14\% attribute steering accuracy (exceeding random chance baselines by 45 percentage points). Linear probing revealed 100\% accuracy in classifying preferences from hidden states, demonstrating that joint training creates shared representations encoding task-relevant information.

We make three primary contributions. First, we provide the first empirical demonstration that DPO preference optimization and attribute conditioning can be jointly trained without catastrophic interference, quantified through gradient angles averaging 78.5° across 100 proof-of-concept steps. Second, we validate that gradient compatibility (<120° angle) serves as a quantitative design principle for predicting multi-objective LLM alignment feasibility, providing a transferable measurement methodology applicable beyond the specific DPO-attribute combination tested. Third, we establish a proof-of-concept methodology for LLM hypothesis validation where 100-step experiments validate feasibility (convergence, gradient compatibility) while deferring performance optimization (achieving baseline parity) to full-scale training.

Our findings are necessarily limited by proof-of-concept scale constraints. Training at 100 steps (approximately 1\% of the planned 15,000-step full training) prevents us from claiming performance parity with standalone baselines or demonstrating emergent benefits over sequential training. The gradient compatibility measurement---architecture-agnostic and independent of training duration---remains our most robust finding, providing strong evidence that full-scale joint training is viable given sufficient computational resources.
